\documentclass[11pt,a4paper]{article}
\usepackage[margin=23mm]{geometry}
\usepackage{amsmath,amssymb,booktabs,microtype}
\usepackage[hidelinks]{hyperref}
\newcommand{\id}{\mathbb I}
\newcommand{\tr}{\operatorname{tr}}
\newcommand{\ket}[1]{|#1\rangle}
\newcommand{\bra}[1]{\langle#1|}
\newcommand{\pr}[1]{\ket{#1}\bra{#1}}
\title{Route G-R2: Two Events, Two Durable Records}
\author{A finite relational protocol, not a density--time law}
\date{24 September 2026}
\begin{document}
\maketitle
\begin{abstract}
A cyclic clock reading is not itself a saved event. We construct one
initial record and one terminal readout, including their quantum
backreaction, inside an exact finite history constraint. Orthogonal
memory sectors retain the outcomes and two fixed protocol stamps.
A one-bit epoch label separates repeated clock phases. The construction
is conditional on an ordered protocol and a blank-memory boundary; it
does not derive an arrow of time, apparatus energetics, or an $E_d$
law. In an exact three-level example the terminal probability is $0$
without the initial record, $2/9$ with it, and $4/9$ for a different
instrument having the same initial measurement effects.
\end{abstract}

\section{Declared system, clock and recording gates}
G-R1 supplied a noninteracting finite Page--Wootters constraint and
single-reading statistics. Recording requires an additional coupling.
Here we retain its three-level free conditional propagator,
\begin{equation}
 H_S=\epsilon\operatorname{diag}(0,1,2),\qquad
 U=\operatorname{diag}(1,z,z^2),\qquad z=e^{-2\pi i/3},\quad U^3=\id ,
\end{equation}
but \emph{replace} its constraint by the explicit protocol constraint
below. No claim is made that the old $H_C+H_S-E_*$ implements these gates.
All phases are clock labels, not spatial rotations.

The work space is $W=S\otimes A\otimes F$, of dimension $27$.
Each memory has orthogonal states $\ket{\emptyset},\ket0,\ket1$.
Let $S_a$ exchange $\ket{\emptyset}$ and $\ket a$, leaving the third
state fixed. With $\ket+=(\ket0+\ket1+\ket2)/\sqrt3$, set
\begin{align}
 P_0&=\pr0,& P_1&=\id-P_0,&
 Q_+&=\pr+,& Q_-&=\id-Q_+ ,\\
 W_i&=\sum_{a=0,1}P_a\otimes S_a\otimes\id,&
 W_f&=\sum_{b=+,-}Q_b\otimes\id\otimes S_b .
\end{align}
For $F$, $S_+$ denotes $S_0$, and $S_-$ denotes $S_1$.
Orthogonality and completeness of $P_a,Q_b$ prove
$W_i^\dagger W_i=W_f^\dagger W_f=\id$. These are unitary
premeasurements; a later pointer read gives the outcome.
For a blank memory,
\begin{equation}
 W_i\ket\psi\ket\emptyset_A\ket\emptyset_F
 =\sum_a P_a\ket\psi\ket a_A\ket\emptyset_F .
\end{equation}
This is standard instrument physics, not a new measurement law.
Explicit memories in conditional-time models are discussed in
\cite{quantumtime}; our finite boundary/epoch choice is declared here.

\newpage
\section{An exact finite stationary history, with an explicit boundary}
Take the orthonormal G-R1 clock grid
$\ket{\varphi_k}=3^{-1/2}\sum_{n=0}^2e^{2\pi i nk/3}\ket n$.
Add an epoch bit $R$, and use
\begin{equation}
 \ket q=\ket{\varphi_{q\bmod3}}_C
       \ket{\lfloor q/3\rfloor}_R,\qquad 0\leq q\leq5.
\end{equation}
This is six distinct subsystem configurations, not six externally
calibrated times. The bit is not an unlimited cycle counter.
For $\chi_0=\ket+\ket\emptyset_A\ket\emptyset_F$ define
$U_W=U\otimes\id_A\otimes\id_F$ and
\begin{equation}\label{eq:links}
 (V_0,V_1,V_2,V_3,V_4)=(W_i,U_W,W_fU_W,U_W,U_W),
 \qquad \chi_{q+1}=V_q\chi_q .
\end{equation}
Only one initial gate and one final gate are applied. Two free
propagators separate the recorded initial state $\chi_1$ from
the terminal measurement in $\chi_3$; hence $U_\Delta=U^2$.
Recording links are not asserted to have the duration of a free tick.

The memory labels may be read as $\ket{(r_i,k_i),a}_A$ and
$\ket{(r_f,k_f),b}_F$, with fixed stamps $(r_i,k_i)=(0,1)$ and
$(r_f,k_f)=(1,0)$, using the post-gate configuration convention.
Because each event has a predeclared stamp, three memory states per
event suffice. This is \emph{not} an autonomous recorder that measures
an arbitrary unknown clock reading. The gate position supplies the
stamp; a general timestamp recorder would require additional states.

Define row maps and a positive operator of dimension $162$:
\begin{align}
 A_q&=\bra{q+1}\otimes\id-\bra q\otimes V_q,\\
 K_{\rm prop}&=\frac12\sum_{q=0}^4 A_q^\dagger A_q,\qquad
 K_{\rm in}=\pr0\otimes(\id-\pr{\chi_0}),\\
 K_{\rm hist}&=K_{\rm prop}+K_{\rm in},\qquad
 \ket\Omega=\frac1{\sqrt6}\sum_{q=0}^5\ket q\chi_q .
\end{align}
This is a finite history-state construction in the tradition of
\cite{history}, with all terms specified rather than a continuum ideal
clock. For any $\ket X=\sum_q\ket q x_q$,
\begin{equation}\label{eq:squares}
 \bra X K_{\rm hist}\ket X
 =\frac12\sum_{q=0}^4\|x_{q+1}-V_qx_q\|^2
  +\|(\id-\pr{\chi_0})x_0\|^2\geq0 .
\end{equation}
Thus $K_{\rm hist}\Omega=0$. Conversely, zero energy in
\eqref{eq:squares} implies $x_0=c\chi_0$ and $x_{q+1}=V_qx_q$,
so the kernel is exactly one-dimensional.
Conditioning on $q$ gives $\chi_q$ and $p(q)=1/6$.

The ordered links and blank-input boundary are assumptions, not an
emergent causal order. Positivity of this protocol constraint is not
an apparatus energy budget, and its eigenvalues have no assigned
physical energy scale. We have not retained G-R1's separate proof
of a noninteracting, fixed-total-local-energy realization.

\newpage
\section{Sequential probabilities and an exact backreaction test}
For initial density matrix $\rho$, instruments with Kraus operators
$M_a,N_b$ satisfying $\sum_aM_a^\dagger M_a=\sum_bN_b^\dagger N_b=\id$
give the joint memory statistics
\begin{equation}\label{eq:joint}
 p(a,b)=\tr\!\left[
 N_bU_\Delta M_a\rho M_a^\dagger U_\Delta^\dagger N_b^\dagger
 \right] .
\end{equation}
Indeed, for pure input the completed record state is
$\sum_{a,b}N_bU_\Delta M_a\ket\psi\ket a_A\ket b_F$.
The squared norm in the orthogonal pointer sector $(a,b)$ yields
\eqref{eq:joint}; linearity gives mixed inputs. Completeness gives
\begin{equation}
 \sum_bp(a,b)=\tr(M_a\rho M_a^\dagger)=p(a),\quad
 \sum_{a,b}p(a,b)=1,\quad
 p(b\mid a)=p(a,b)/p(a)
\end{equation}
when $p(a)>0$. These are standard sequential Born probabilities
\cite{quantumtime}, not a product of separate unperturbed readings.
Here they also follow directly by projecting $\Omega$ onto $q=3$
and the two pointer sectors.

For our projective instrument $M_a=P_a,N_b=Q_b$ and $\rho=\pr+$,
\begin{align}
 \langle+|U^2P_0|+\rangle&=1/3,\\
 \langle+|U^2P_1|+\rangle&=(z^2+z^4)/3=-1/3 .
\end{align}
Consequently
\begin{equation}\label{eq:table}
 [p(a,b)]_{\substack{a=0,1\\b=+,-}}
 =\frac19\begin{pmatrix}1&2\\1&5\end{pmatrix},\qquad
 p(a)=(1/3,2/3),\quad
 p(+\mid a)=(1/3,1/6).
\end{equation}
The minus entries are the corresponding branch norms
$\|P_a\ket+\|^2$ minus the plus entries. Ignoring the first
pointer, but actually creating its record, gives
\begin{equation}
 \rho\longmapsto\sum_aP_a\rho P_a,\qquad p(+)=2/9 .
\end{equation}
By contrast, omitting the initial interaction altogether gives
\begin{equation}
 p_{\rm no\ record}(+)=|\langle+|U^2|+\rangle|^2
 =|1+z^2+z^4|^2/9=0 .
\end{equation}
The difference is measurement backreaction, not a changed mass or a
slower clock. The rank-two branch retains coherence internally;
the initial instrument removes only coherence between the two
$P_a$ sectors.

\subsection*{Why the effects alone do not specify the protocol}
Take $R_0=\id$, $R_1=\operatorname{diag}(1,1,-1)$ and
$M'_a=R_aP_a$. Then $M_a'{}^\dagger M'_a=P_a$:
the initial outcome probabilities are unchanged. But the plus
amplitude in branch $1$ is $(z^2-z^4)/3$, of squared norm $1/3$.
Thus
\begin{equation}
 [p'(a,b)]=\frac19\begin{pmatrix}1&2\\3&3\end{pmatrix},
 \qquad p'(+)=4/9 .
\end{equation}
A controlled $R_a$ on the already written pointer implements this
alternative unitarily. It is an audit of instrument dependence, not
an activated $E_d$ interaction.

\newpage
\section{Saved records, recurrent readings and the cycle obstruction}
Let $\Pi_a^A=\id_S\otimes\pr a_A\otimes\id_F$ and define $\Pi_b^F$
similarly, including the blank sector. Every link after writing $A$
commutes with $\Pi_a^A$; every link after writing $F$ commutes with
both pointers. Hence
\begin{equation}
 \langle\chi_q|\Pi_a^A\Pi_b^F|\chi_q\rangle=p(a,b),
 \qquad q=3,4,5.
\end{equation}
In fact the entire reduced state of $A$ is unchanged by any subsequent
unitary acting only on $S,F$. Conditioning on a terminal outcome can
update our inferred distribution of $a$ but does not overwrite $A$.
Repeated projective pointer reads are nondemolition since
$\Pi_a\Pi_{a'}=\delta_{aa'}\Pi_a$.
The nonselective read map
$\rho_W\mapsto\sum_{a,b}\Pi_a^A\Pi_b^F\rho_W\Pi_a^A\Pi_b^F$
preserves all pointer probabilities.

This does not assert that the joint memory density matrix has already
lost every off-diagonal entry: premeasurement can retain coherences.
Nor is an arbitrary quantum state cloned. The records are distinguishable
pointer sectors with a repeatable classical readout.

With $\Pi_{\rm filled}=\id-\pr\emptyset$, define the two-event counter
\begin{equation}
 N_{\rm evt}=\id_S\otimes\Pi_{\rm filled}^A\otimes\id_F
            +\id_S\otimes\id_A\otimes\Pi_{\rm filled}^F .
\end{equation}
Each slice is an exact eigenstate, $N_{\rm evt}\chi_q=n_q\chi_q$:
\begin{center}
\begin{tabular}{c*{6}{c}}
\toprule
configuration $q$&0&1&2&3&4&5\\
clock reading $k$&0&1&2&0&1&2\\
epoch $r$&0&0&0&1&1&1\\
stored-event count $n_q$&0&1&1&2&2&2\\
\bottomrule
\end{tabular}
\end{center}
Only the two declared recording events are counted, not every phase
tick. This is not a universal identification with the source's $T$
or energy information.

Ignoring the epoch and conditioning only on $k=0$ leaves
\begin{equation}
 \rho_{W\mid k=0}=\tfrac12
     (\pr{\chi_0}+\pr{\chi_3}),\qquad
 \Pr(F\ {\rm filled}\mid k=0)=\tfrac12 .
\end{equation}
Reading $r=1$ as well, or selecting the completed terminal-memory
sector, instead singles out $\chi_3$ at this phase. A cyclic phase
without a record or a cycle tag does not identify a unique event.

Suppose instead we forced the entire work state to be periodic after
this episode while starting both memories blank. Closure would require
\begin{equation}
 W_{\rm cyc}\chi_0=\chi_0,\qquad W_{\rm cyc}=W_fU_W^2W_i .
\end{equation}
But the left side has both pointers filled, whereas the right side has
both blank. They are orthogonal, so
\begin{equation}
 \langle\chi_0|W_{\rm cyc}\chi_0\rangle=0,\qquad
 \|(W_{\rm cyc}-\id)\chi_0\|=\sqrt2 .
\end{equation}
The incompatible requirements are full-state recurrence and a fresh,
retained record, not cyclic clocks in general. Resetting, unwriting,
a longer recurrence, or an explicit epoch/memory extension changes
the requirements. Our open history uses the last option.

\newpage
\section{What this closes, and what it deliberately leaves open}
The bounded G-R2 question is answered: specified record interactions
convert two scheduled cyclic-clock events into orthogonal, readable
memory sectors. One static joint history yields the complete
two-event Born statistics, including backreaction and retained records.
The physical information added to a recurrent phase is
\emph{which episode, which event, and which instrument outcome},
not a fitted function of energy density.

There are three important qualifications.
\begin{enumerate}
 \item \emph{Finite protocol, not an autonomous universal clock.}
 The epoch bit, two fixed stamps, ordered gates and input boundary
 are explicit inputs. No general unknown timestamp recorder, infinite
 counter or uniquely derived event algebra is claimed. We do not
 claim the separate epoch register is dimension-minimal; in other
 protocols memory occupancy can also distinguish stages.
 \item \emph{Preservation, not thermodynamic irreversibility.}
 The records persist under the specified future links and could
 persist under further isolated-pointer links. Real memory lifetime,
 environmental errors, redundant records and erasure costs are not
 supplied. In this model $W_{\rm cyc}^\dagger$ restores the blank
 state exactly. The stationary history and its oriented boundary
 therefore do not derive a thermodynamic arrow.
 \item \emph{Constraint consistency, not physical energy completion.}
 $[W_i,H_S\otimes\id_{AF}]=0$, but
 $[W_f,H_S\otimes\id_{AF}]\neq0$ because
 $[Q_+,H_S]\neq0$. Work/clock/apparatus degrees of freedom are needed
 for a physical energy-conserving realization of that readout.
 $K_{\rm hist}\Omega=0$ does not by itself supply this realization.
 A microscopic clock speed or gate duration cannot be extracted
 from its arbitrary overall normalization.
\end{enumerate}

\subsection*{Verification and reproducibility}
The executable companion
\texttt{route\_g/code/verify\_gr2\_two\_event\_records.py}
constructs the gates and the full $162\times162$ constraint independently
of the analytic probabilities. Its 139 checks cover unitarity,
projector completeness, sum-of-squares assembly, the unique zero mode,
all conditioned slices, both exact probability tables, the no-record
test, memory preservation, phase aliasing, failed periodic closure,
reversal, mixed inputs and unchanged prior source files.
These certify finite algebra, not empirical physics.
The JSON report retains individual residuals and source hashes.

\subsection*{Next decision, not executed in G-R2}
For G-R3, choose an explicit $E_i/E_d$ object and reference measure,
then one clock interaction and a reference-clock comparison.
In particular, distinguish a state-dependent interaction from a common
energy rescaling or a relabeling of one clock. No slowdown law is
selected, fitted or inferred here. The original density--time proposal,
emergent spacetime, gravity, and all Route-F promotion gates remain open.

\begin{thebibliography}{9}
\bibitem{quantumtime}
V.~Giovannetti, S.~Lloyd and L.~Maccone,
\emph{Quantum Time}, Phys.\ Rev.\ D \textbf{92}, 045033 (2015),
\href{https://arxiv.org/abs/1504.04215v3}{arXiv:1504.04215v3},
especially the Measurements section, Eqs.~(31)--(46).
\bibitem{history}
D.~Aharonov et al.,
\emph{Adiabatic Quantum Computation is Equivalent to Standard
Quantum Computation}, SIAM J.\ Comput.\ \textbf{37}, 166--194 (2007),
\href{https://arxiv.org/abs/quant-ph/0405098}{arXiv:quant-ph/0405098}.
\end{thebibliography}
\end{document}
